\chapter*{Introduction g\'en\'erale}
\addcontentsline{toc}{chapter}{Introduction g\'en\'erale}
\markboth{Introduction g\'en\'erale}{Introduction g\'en\'erale}

% ---------------------------------------------------------------------
% GUIDE DE REDACTION (a supprimer une fois le texte ecrit)
% L'introduction generale (1 a 2 pages) doit repondre a 4 questions :
%   1. Contexte : dans quel domaine se situe le projet ?
%      (commerce international, douanes, digitalisation des processus)
%   2. Probleme : quelle difficulte rencontrent les acteurs aujourd'hui ?
%      (gestion manuelle des regles douanieres, processus non traces,
%       creation de regles necessitant des competences techniques DRL)
%   3. Solution : que propose le projet E-FORCE ?
%      (plateforme web centralisee : moteur de regles metier + BPMN
%       + module IA/NLP + import/export de dossiers)
%   4. Plan : annoncer les chapitres du rapport.
% ---------------------------------------------------------------------

Le commerce international repose sur des proc\'edures douani\`eres
complexes, encadr\'ees par de nombreuses r\`egles m\'etier (contr\^oles,
taxation, certification, quotas) ainsi que par des processus
organisationnels impliquant plusieurs acteurs. [\`A compl\'eter : situer le
contexte g\'en\'eral du projet et de l'organisme d'accueil.]

[\`A compl\'eter : d\'ecrire la probl\'ematique -- par exemple la difficult\'e
de maintenir et de faire \'evoluer des r\`egles m\'etier exprim\'ees dans un
langage technique (DRL/Drools), l'absence de visualisation claire des
processus m\'etier, et la gestion dispers\'ee des dossiers
d'import/export.]

Pour r\'epondre \`a ces besoins, ce projet propose la conception et le
d\'eveloppement de la plateforme \textbf{E-FORCE}, une application web
permettant de centraliser la gestion des processus m\'etier (mod\'elis\'es
en BPMN), la gestion des r\`egles m\'etier (avec g\'en\'eration automatique de
code DRL et un module d'intelligence artificielle de traitement du
langage naturel), ainsi que la gestion des dossiers d'importation et
d'exportation.

Ce rapport est organis\'e en quatre chapitres. Le premier chapitre
pr\'esente le contexte g\'en\'eral du projet, l'organisme d'accueil, l'\'etude
de l'existant et la solution propos\'ee. Le deuxi\`eme chapitre est
consacr\'e \`a l'analyse et \`a la sp\'ecification des besoins fonctionnels et
non fonctionnels. Le troisi\`eme chapitre d\'ecrit la conception de la
solution (architecture, mod\`ele de donn\'ees, diagrammes UML). Enfin, le
quatri\`eme chapitre pr\'esente la r\'ealisation, les technologies utilis\'ees
et les principales fonctionnalit\'es d\'evelopp\'ees, illustr\'ees par des
captures d'\'ecran.

\clearpage
